\documentclass[11pt]{article}

\usepackage[margin=1in]{geometry}
\usepackage{listings}
\usepackage{xcolor}
\usepackage{hyperref}
\usepackage{amsmath}
\usepackage{parskip}

\hypersetup{
  colorlinks=true,
  linkcolor=blue!60!black,
  urlcolor=blue!60!black
}

\lstset{
  basicstyle=\ttfamily\small,
  frame=single,
  breaklines=true,
  showstringspaces=false,
  columns=fullflexible,
  keywordstyle=\color{blue!70!black},
  commentstyle=\color{green!40!black},
  stringstyle=\color{red!60!black}
}

\title{CS201 Spring 2026 --- Lecture Notes}
\author{Clarissa Littler}
\date{Spring 2026}

\begin{document}
\maketitle

\section*{Preface}

This document is a running narrative of what we covered in each class
meeting this term. Each section pairs the agenda for the day with the
little code samples we actually typed up together, so if you missed a
lecture or want to go back and re-read, the whole arc of a day should be
here in one place. The source files themselves live in the
\texttt{lecture$N$/} subdirectories of the class repo, and all of the
snippets quoted below are verbatim from those files.

\tableofcontents

\section{Lecture 1 --- Welcome, C vs.\ C++, Bitwise Tricks, and a Peek at Integers}

This one was mostly an orientation plus a crash-course in the parts of
C that tend to trip up people coming from C++. I was genuinely a little
surprised everyone showed up, so thank you for that!

\subsection*{Logistics}

The class is my own personal admixture of four things:
\begin{itemize}
  \item C
  \item x86-64 assembly
  \item Concurrency
  \item Operating systems theory and practice
\end{itemize}
Grading-wise, you've got quizzes, discussions, and coding assignments.
The class repo is going to be your friend: that's where the textbook
lives, where all the ``Agenda'' files like this one come from, and
where we'll drop every bit of code we work on in class.

Structurally, I'll aim to start about five minutes after 10, break
after every 50 minutes, and try not to keep you much past two hours.

\subsection*{C: What's different?}

If you're coming from C++, the big shocks are:
\begin{itemize}
  \item No pass-by-reference. You get pass-by-value or you get
    pointers, full stop.
  \item No \texttt{bool} type (well, not in the way you're used to ---
    nonzero is truthy and zero is falsy, and there's no native
    \texttt{true}/\texttt{false}).
  \item No \texttt{nullptr}; you use the \texttt{NULL} macro (which is
    basically just \texttt{0} cast to a pointer).
  \item I/O is \texttt{printf}/\texttt{scanf} rather than streams.
\end{itemize}

We played with two small files to make this concrete. The first,
\texttt{arraything.c}, shows that arrays do decay to pointers --- so
you pass the length by hand --- and that you iterate with a plain
\texttt{for} loop rather than an iterator:

\begin{lstlisting}[language=C]
void arrayMunger(int a[], int s){
  for(int i=0; i<s; i++){
    a[i] = 2*a[i];
  }
}
\end{lstlisting}

The little commentary \texttt{// arr[i] ===> *(arr + i)} hints at the
thing we'll lean on all term: indexing is just pointer arithmetic in
disguise.

The second file, \texttt{passPointy.c}, makes the ``no pass by
reference'' point tangible. If you want a function to modify its
caller's variable, you pass the \emph{address} and dereference inside:

\begin{lstlisting}[language=C]
void intMunger(int* p){
  *p = 2*(*p);
}
\end{lstlisting}

And you call it with \texttt{intMunger(\&test);}. That ampersand is
doing real work --- it's giving you a pointer to \texttt{test} rather
than a copy of it.

\subsection*{Bitwise operations}

We went through the usual suspects with truth tables:
\begin{itemize}
  \item \texttt{\&} (AND), \texttt{|} (OR), \texttt{\^{}} (XOR),
    \texttt{\textasciitilde{}} (NOT).
  \item \texttt{<<} (shift left, pads with zeroes, anything that
    falls off the edge is gone) and \texttt{>>} (shift right, same
    idea on the other side).
\end{itemize}

The file \texttt{bitmunger1.c} is where this all came together. It
defines a little loop that prints out the 32 bits of an \texttt{int},
lets you pick a bit to toggle, and optionally negates the integer via
two's complement:

\begin{lstlisting}[language=C]
void toggleBit(int* num,int place){
  *num = (*num) ^ (1 << place);
}

void twosComplement(int* num){
    *num = ~((*num) - 1);
}
\end{lstlisting}

\texttt{toggleBit} is the canonical XOR trick: XOR-ing a bit with 1
flips it, XOR-ing with 0 leaves it alone. \texttt{twosComplement}
implements the ``subtract one, then invert'' rule for negating a
two's-complement integer.

\subsection*{Peeking inside an integer: two's complement}

We derived why two's complement works. A plain unsigned $n$-bit
integer encodes
\[
  \sum_{i=0}^{n-1} b_i \cdot 2^i,
\]
so \texttt{11111111} (8 bits) is $128 + 64 + \cdots + 1 = 255$. In
two's complement, the top bit's weight flips sign:
\[
  -b_{n-1} \cdot 2^{n-1} + \sum_{i=0}^{n-2} b_i \cdot 2^i.
\]
That gives you one clean mental check: all-ones is always $-1$,
because $-2^{n-1} + (2^{n-1} - 1) = -1$. And to negate a number, you
subtract one then invert all the bits (or equivalently, invert and
then add one).

\subsection*{A preview of floats}

We ran out of time before we could fully explore floats, but I showed
off \texttt{floatmunger1.c}, which is structurally the same as
\texttt{bitmunger1.c} but reinterprets an \texttt{int}'s bits as a
\texttt{float} via a pointer cast. The little spaces it prints
(between bits 31/30 and 23/22) are hinting at the IEEE-754 sign /
exponent / fraction split we'll dig into next time.

\section{Lecture 2 --- Hex, Endianness, IEEE-754, and Where Floats Lie to You}

The plan was to finish the float discussion we previewed on day one,
but before diving in we took a small detour through hex and
endianness.

\subsection*{Warm-up: hex codes and the shape of a pointer}

Hex is base-16, digits 0--F. The reason it shows up everywhere ---
addresses, colors, byte dumps --- is that two hex digits are exactly
one byte ($16 \cdot 16 = 256 = 2^8$). So when you see an RGBA color
like \texttt{\#8765AF}, that's really three (or four) bytes: one for
red, one for green, one for blue, optionally one for transparency.

\subsection*{Endianness}

We wrote \texttt{byteOrderTest.c} to actually look at how an 8-byte
integer is laid out in memory on our machines:

\begin{lstlisting}[language=C]
long int num = 0x0123456789abcdef;
unsigned char* p = (unsigned char*)&num;

printf("The bytes are stored in memory as: ");
for(int i=0; i<8; i++){
  printf("%x ",*(p+i));
}
\end{lstlisting}

On an x86-64 box this prints the bytes out in the opposite order from
how we wrote the literal --- classic little-endian. The trick is the
\texttt{(unsigned char*)} cast: once we have a byte-pointer, pointer
arithmetic moves us one byte at a time, so we can walk the actual
storage.

\subsection*{Sizes and pointers}

Alongside that, \texttt{sizes.c} just prints \texttt{sizeof(int)},
\texttt{sizeof(long)}, \texttt{sizeof(char)}, \texttt{sizeof(float)},
\texttt{sizeof(double)}, and the size of a pointer. On our 64-bit
Linux boxes: int is 4, long is 8, char is 1, float is 4, double is 8,
and any pointer is 8. That last one is the important bit ---
\emph{all} pointers are the same size, regardless of what they point
to, because they're all just addresses.

And \texttt{pointerArith.c} drove home that pointer arithmetic scales
with the \emph{pointee} size. An \texttt{int*} advances 4 bytes per
\texttt{+1}, a \texttt{char*} advances 1, a \texttt{char**} advances
8, and so on.

\subsection*{IEEE-754, formally}

We did the 32-bit (\texttt{float}) version because the 64-bit version
is just ``same thing with bigger fields.''

\begin{itemize}
  \item 1 bit sign
  \item 8 bits exponent (stored with a bias of 127)
  \item 23 bits fraction
\end{itemize}

In the normal case, the value is
\[
  (-1)^s \cdot 2^{e - 127} \cdot (1 + f),
\]
where $f$ is the fractional part read as ``rational binary'': the bit
at position $k$ after the implicit point contributes $2^{-k}$.

Worked example: \texttt{0 10000001 1010\ldots0} is
\[
  (-1)^0 \cdot 2^{129 - 127} \cdot (1 + \tfrac12 + \tfrac18)
  = 4 \cdot \tfrac{13}{8} = 6.5.
\]

We exercised this live with \texttt{floatmunger1.c} (same idea as
lecture 1, but now with \texttt{\%.16f} so you can see lots of digits).
Flipping specific bits and watching the printed value change is the
single best way to internalize the format.

\subsubsection*{Edge cases}

\begin{itemize}
  \item Exponent is all 1s:
    \begin{itemize}
      \item fraction is zero $\Rightarrow$ $\pm\infty$,
      \item fraction is nonzero $\Rightarrow$ \texttt{NaN}.
    \end{itemize}
  \item Exponent is all 0s (subnormal numbers): the exponent sticks at
    $2^{-126}$, but the leading implicit 1 is gone, so the value is
    $(-1)^s \cdot 2^{-126} \cdot f$.
\end{itemize}

\subsection*{Where floats lie to you}

\texttt{floatlimit.c} is the ``gotcha'' example:

\begin{lstlisting}[language=C]
float f = 0.0001;
float accum = 0;
for(int i=0; i<10000; i++){
  accum += f;
}
printf("%.20f\n",accum);
\end{lstlisting}

We'd love this to print \texttt{1.00000000000000000000}. It doesn't.
The reason is a mix of: 0.0001 is not exactly representable in binary,
and we're accumulating that small rounding error 10{,}000 times.

Why is this \emph{unavoidable}? A bit of math: the rationals and the
integers are countably infinite, but the reals form a strictly larger
uncountable set. There is no way to cover the reals with finitely many
bits, or even with countably many bits; we're forced to pick a
countable scatter of ``breadcrumbs'' and round every other real to the
nearest one.

\subsubsection*{Round to even}

To keep that rounding error from biasing one direction, IEEE-754
specifies ``round to even'': when a real is exactly halfway between
two breadcrumbs, the hardware rounds to the one whose last bit is
zero. Over many operations this rounds up about as often as it rounds
down.

\section{Lecture 3 --- \texttt{malloc}, and Your First Real Assembly}

The real headline today: we left C-land and wrote actual assembly
against the bare Linux syscall interface. But first, I owed you
\texttt{malloc}.

\subsection*{\texttt{malloc}, finally}

In my rush last time I completely skipped dynamic allocation, which is
embarrassing, so we opened with two quick examples.

\texttt{malloc.c} is the minimum viable pointer dance:

\begin{lstlisting}[language=C]
int* n = malloc(sizeof(int));
*n = 10;
printf("%d\n",*n);
free(n);
n = 0;
if(0 == NULL){
  printf("Yup\n");
}
\end{lstlisting}

Two things to note. First, \texttt{malloc} always returns a pointer
and always needs you to tell it how many bytes you want --- unlike
\texttt{new T} in C++, it doesn't know anything about your type, so
you pass \texttt{sizeof(int)} (or a multiple of it). Second, that
final \texttt{if} is my excuse to point out that \texttt{NULL} is
literally just \texttt{0} cast to a pointer type.

\texttt{malloc2.c} extends this to an array and a struct:

\begin{lstlisting}[language=C]
int* narr = malloc(100*sizeof(int));
narr[2] = 5;

struct Goofus* g = malloc(sizeof(struct Goofus));
g->g1 = 5;            // same thing as (*g).g1
\end{lstlisting}

Once you've got a pointer from \texttt{malloc}, you can index it just
like an array. The \texttt{->} arrow is shorthand for ``dereference
then access a field.''

\subsection*{Stepping back: what is a CPU, anyway?}

The big conceptual pivot for the rest of the term:
\begin{itemize}
  \item A CPU is, at bottom, a thing that fetches \emph{instructions}
    from memory and does what they say.
  \item The set of instructions it understands is an \emph{instruction
    set architecture} (ISA). Ours is x86-64.
  \item Assembly language is a one-to-one (ish) textual encoding of
    those instructions, so writing assembly is as close to ``talking
    to the CPU'' as we'll practically get.
  \item Registers are a tiny amount of extremely fast storage inside
    the CPU itself. Most instructions operate on registers; memory
    comes in via load/store.
  \item Memory hierarchy, the 30-second version: registers $\to$
    L1/L2/L3 cache $\to$ main memory $\to$ disk. Each level is
    roughly an order of magnitude slower and roughly an order of
    magnitude bigger than the one above it.
\end{itemize}

I also mentioned \texttt{gcc -S}, which tells gcc to stop at assembly
rather than assemble and link. We'll come back to it.

\subsection*{Your first bare assembly program}

\texttt{first.s} is deliberately broken:

\begin{lstlisting}
    .section .text
    .global _start

_start:
    mov $60,%rax
\end{lstlisting}

This is the shortest thing that looks like an assembly program. The
\texttt{.text} section is where code lives; \texttt{\_start} is the
symbol the linker uses as the program entry point (unlike in a C
program, there is no runtime to call \texttt{main} for us).

We moved \texttt{60} into \texttt{\%rax} --- 60 is the syscall number
for \texttt{exit} on x86-64 Linux --- and then\ldots{} did nothing
else. Because we never actually issued the \texttt{syscall}
instruction and there's no implicit return, the CPU just wandered off
the end of our code and trapped. Hence the crash.

\texttt{firstfixed.s} is the corrected version:

\begin{lstlisting}
_start:
    mov $60,%rax
    mov $-1,%rdi
    syscall
\end{lstlisting}

Syscall convention: put the syscall number in \texttt{\%rax}, put
arguments in \texttt{\%rdi}, \texttt{\%rsi}, \texttt{\%rdx}, etc., and
then execute \texttt{syscall}. For \texttt{exit(status)}, the status
goes in \texttt{\%rdi}. So this exits with status \texttt{-1}, and
\texttt{echo \$?} in the shell will confirm it.

\subsection*{\texttt{add} and \texttt{sub}}

\texttt{firstadd.s} and \texttt{firstsub.s} are our first binary
operations:

\begin{lstlisting}
_start:
    mov $10,%rbx
    mov $20,%rcx
    add %rbx,%rcx
    mov %rcx,%rdi
    mov $60,%rax
    syscall
\end{lstlisting}

Syntax note that will bite you all term: in AT\&T syntax,
\texttt{add S, D} means \texttt{D = D + S}. The destination is on the
\emph{right}. So \texttt{add \%rbx,\%rcx} leaves 30 in
\texttt{\%rcx}. \texttt{sub} is analogous: \texttt{sub \%rbx,\%rcx}
computes \texttt{D - S} and leaves 10 in \texttt{\%rcx}.

We also did the fun thing of running \texttt{gcc -S test.c} on a tiny
C file and skimming the output (\texttt{test.s}). It's noisy
(\texttt{.cfi\_}directives, frame setup, PLT calls) but the actual
work is recognizable: move 10 into a local, load it into
\texttt{\%esi}, load the format string's address via \texttt{leaq},
call \texttt{printf@PLT}.

\section{Lecture 4 --- Registers, Calling Convention, Addressing Modes, Control Flow, and ``Hello, world''}

Today was a dense one: we filled in the register table, introduced the
Linux x86-64 calling convention, walked through every addressing mode
we'll use for arrays, covered \texttt{cmp}/\texttt{jmp}, wrote a
for-loop in assembly, and finished by doing \texttt{write} and
\texttt{exit} syscalls by hand.

\subsection*{Logistics}

The Linux server admin is in the process of activating accounts for
folks who can't log in yet, so hang tight if that's you.

\subsection*{Register reference}

The basic x86-64 general-purpose registers and what they were
historically used for:

\begin{itemize}
  \item \texttt{\%rax} --- accumulator (return values, implicit operand
    for \texttt{mul}/\texttt{div}).
  \item \texttt{\%rbx} --- base register (historically an array base).
  \item \texttt{\%rcx} --- counter (loops, \texttt{rep}).
  \item \texttt{\%rdx} --- data (I/O, upper half of
    multiplication/division).
  \item \texttt{\%rsi} --- source index (string ops).
  \item \texttt{\%rdi} --- destination index (string ops).
  \item \texttt{\%rbp} --- base pointer (stack frame base).
  \item \texttt{\%rsp} --- stack pointer (top of the stack).
  \item \texttt{\%r8}--\texttt{\%r15} --- extra general-purpose
    registers.
  \item \texttt{\%rip} --- instruction pointer; we don't write to this
    directly, but we \emph{do} use it in RIP-relative addressing.
\end{itemize}

The floating-point/SIMD registers (\texttt{\%xmm0}\ldots{}) live in
the \texttt{assemblyGuide.org} in the repo.

\subsection*{Linux/x86-64 calling convention, first pass}

\begin{itemize}
  \item Return values go in \texttt{\%rax}.
  \item The first six integer/pointer arguments are passed in, in
    order:
    \texttt{\%rdi}, \texttt{\%rsi}, \texttt{\%rdx}, \texttt{\%rcx},
    \texttt{\%r8}, \texttt{\%r9}.
\end{itemize}

Syscalls use a slightly different convention (the fourth argument is
\texttt{\%r10}, not \texttt{\%rcx}), but for function calls in user
code the list above is what you want.

\subsection*{\texttt{add}/\texttt{sub}, again}

We re-ran \texttt{firstadd.s} from last time as a warm-up; the only
new thing is a pair of clarifying comments explaining that
\texttt{\%rdi} is the exit-code argument and \texttt{60} in
\texttt{\%rax} selects the \texttt{exit} syscall.

\subsection*{Global variables: values, pointers, and \texttt{lea}}

Three files building on each other:

\texttt{addGlobal.s} --- define a global with \texttt{.quad}, then
load, add, store, load-again:

\begin{lstlisting}
    .section .data
num:    .quad 200

_start:
    mov num,%rbx       # load num into %rbx
    add $10,%rbx
    mov %rbx,num       # store back
    mov num,%rdi
    mov $60,%rax
    syscall
\end{lstlisting}

This works, but referencing a symbol like \texttt{num} directly makes
the assembler emit an absolute address, which is awkward for
position-independent code.

\texttt{addGlobalBetter.s} --- same thing, but using
\emph{RIP-relative} addressing:

\begin{lstlisting}
    mov num(%rip),%rbx
    add $10,%rbx
    mov %rbx,num(%rip)
\end{lstlisting}

Here \texttt{num(\%rip)} means ``the address that is
\texttt{num}-minus-the-next-instruction bytes away from the current
instruction pointer.'' The assembler computes that offset for us at
build time. This is the form you want in modern code.

\texttt{addGlobalLea.s} --- now using \texttt{lea} to take
\emph{the address of} \texttt{num}, the way \texttt{\&} does in C:

\begin{lstlisting}
    lea num(%rip),%rbx   # rbx = &num
    addq $10,(%rbx)      # *rbx = *rbx + 10
    mov num(%rip),%rdi
    mov $60,%rax
    syscall
\end{lstlisting}

The comment block at the bottom of the file spells out the C
equivalent:
\begin{lstlisting}[language=C]
int num = 200;
int* nump = &num;
*nump = *nump + 10;
\end{lstlisting}

Parentheses around a register mean ``dereference.'' So \texttt{(\%rbx)}
is ``the 8 bytes of memory starting at the address held in
\texttt{\%rbx}.''

\subsection*{From a scalar to an array}

Four files, one new addressing idea at a time.

\texttt{addArray1.s} --- change \texttt{num} to an array of four
quadwords. The code is unchanged, so it just touches the first element.

\texttt{addArray2.s} --- to get at element 1, add 8 to the pointer
(because one \texttt{.quad} is 8 bytes):

\begin{lstlisting}
    lea num(%rip),%rbx
    addq $8,%rbx
    addq $10,(%rbx)
    movq (%rbx),%rdi
\end{lstlisting}

\texttt{addArray3.s} --- same effect, but using the full ``scale-index''
addressing mode:

\begin{lstlisting}
    lea num(%rip),%rbx
    mov $1,%rcx
    addq $10,(%rbx,%rcx,8)     # *(base + index*8)
    movq (%rbx,%rcx,8),%rdi
\end{lstlisting}

The syntax \texttt{(base, index, scale)} computes
\texttt{base + index*scale}. The scale has to be 1, 2, 4, or 8. This
is exactly how C-style array indexing compiles --- you'll see the same
form over and over.

\texttt{addArray4.s} --- same as \texttt{addArray2}, but with
\texttt{.long} (4-byte) elements instead of \texttt{.quad}, which
forces us to use the \texttt{l} suffix (\texttt{addl},
\texttt{movl}) and the 32-bit register name \texttt{\%edi}.

Note the subtle point: with \texttt{.long}, a step of ``one element''
is 4 bytes, not 8. In the file as written, we advanced the pointer by
8 bytes, which lands us at element \emph{two} of the \texttt{.long}
array, not element one. That kind of size/stride mismatch is a classic
bug I'll keep harping on all term.

\subsection*{\texttt{cmp} and \texttt{jmp}}

\texttt{cmp S, D} does (roughly) \texttt{D - S} and sets flags without
saving the result anywhere. The conditional jumps then read those
flags:

\begin{itemize}
  \item \texttt{jmp} --- unconditional.
  \item \texttt{jl} --- $D < S$. (Yes, \texttt{cmp S, D} then ``if
    less'' compares $D$ to $S$; it still surprises me.)
  \item \texttt{jle} --- $D \le S$.
  \item \texttt{je}, \texttt{jne} --- equal, not equal.
  \item \texttt{jg}, \texttt{jge} --- greater, greater-or-equal.
\end{itemize}

The demo, \texttt{cmp1.s}:

\begin{lstlisting}
    mov $10,%rbx
    mov $20,%rcx
    mov $-1,%rdi
    cmp %rbx,%rcx     # flags set as if computing %rcx - %rbx
    jge greater
    mov $1,%rdi
greater:
    mov $60,%rax
    syscall
\end{lstlisting}

Since $20 \ge 10$, we jump to \texttt{greater} and exit with status
\texttt{-1} (or \texttt{255} as far as the shell is concerned).

\subsection*{A for-loop in assembly}

Two flavors of ``sum the numbers from 1 to 10.''

\texttt{sumLoop.s} counts up:

\begin{lstlisting}
    mov $0,%rbx        # accumulator
    mov $1,%rcx        # counter
loopStart:
    add %rcx,%rbx
    add $1,%rcx
    cmp $10,%rcx
    jle loopStart
\end{lstlisting}

\texttt{sumLoopB.s} counts down:

\begin{lstlisting}
    mov $0,%rbx
    mov $10,%rcx
loopStart:
    add %rcx,%rbx
    sub $1,%rcx
    cmp $0,%rcx
    jg loopStart
\end{lstlisting}

Both leave 55 in \texttt{\%rbx}, which we then move into
\texttt{\%rdi} so we can read it out as the exit status.

\subsection*{Buffers, strings, and \texttt{write}}

The grand finale, \texttt{hello.s}:

\begin{lstlisting}
    .section .data
msg:    .asciz "Hello world!\n"
hellolen = . - msg

    .section .text
    .globl _start

_start:
    mov $1,%rax             # syscall number for write
    mov $1,%rdi             # fd = 1 (stdout)
    lea msg(%rip),%rsi      # buffer address
    mov $hellolen,%rdx      # number of bytes
    syscall

    xor %rdi,%rdi           # exit status 0
    mov $60,%rax            # syscall number for exit
    syscall
\end{lstlisting}

Points of interest:
\begin{itemize}
  \item \texttt{.asciz} lays down the characters followed by a null
    terminator.
  \item \texttt{hellolen = . - msg} is a classic assembly idiom: the
    dot is ``current address,'' so subtracting the start of
    \texttt{msg} gives the length in bytes, including the newline and
    the terminator. We use that as the third argument to \texttt{write}.
  \item \texttt{write} is syscall 1 on Linux x86-64. Arguments are
    \texttt{(fd, buf, count)} in \texttt{\%rdi}, \texttt{\%rsi},
    \texttt{\%rdx}.
  \item \texttt{xor \%rdi,\%rdi} is the idiomatic way to zero a
    register --- one byte shorter than \texttt{mov \$0,\%rdi}.
\end{itemize}

Next time we'll close this out with an \texttt{echo}-style program
that \emph{reads} from stdin and writes back, which means handling
buffers properly.

\section{Lecture 5 --- Hello Again, \texttt{echo}, and the Road from Jumps to Proper Functions}

Today was the big one for putting structure around assembly: we
re-did ``hello, world'' carefully, wrote an \texttt{echo} that reads
and writes a buffer, then walked the whole arc from
\texttt{jmp}-style subroutines to \texttt{call}/\texttt{ret} and
finally to real stack frames with local variables. By the end, we
had a function in assembly that looks and behaves like a function in C.

\subsection*{Hello, world --- this time with feeling}

We re-examined \texttt{hello.s} slowly, focusing on the pieces we
glossed over last time:

\begin{lstlisting}
    .section .data
msg:    .asciz "Hello world!\n"
hellolen = . - msg

    .section .text
    .global _start

_start:
    mov $1,%rax        # write syscall
    mov $1,%rdi        # fd = 1 (stdout)
    lea msg(%rip),%rsi # buffer
    mov $hellolen,%rdx # length
    syscall

    xor %rdi,%rdi
    mov $60,%rax
    syscall
\end{lstlisting}

A few things worth restating:
\begin{itemize}
  \item File descriptors in Unix are just integers indexing into a
    per-process table the kernel maintains. \texttt{0} is stdin,
    \texttt{1} is stdout, \texttt{2} is stderr, and anything you
    \texttt{open} returns the next free slot. \texttt{write} doesn't
    care what kind of thing the fd refers to --- pipe, socket, file,
    terminal --- it all looks the same from up here.
  \item The \texttt{write} syscall \emph{returns} something in
    \texttt{\%rax}: the number of bytes it actually wrote (or a
    negative errno on failure). We didn't use it here, but we will.
\end{itemize}

\texttt{helloRet.s} is the tiny variant that does use it:

\begin{lstlisting}
    mov $1,%rax
    mov $1,%rdi
    lea msg(%rip),%rsi
    mov $hellolen,%rdx
    syscall

    mov %rax,%rdi      # return value of write becomes our exit status
    mov $60,%rax
    syscall
\end{lstlisting}

Run it and then \texttt{echo \$?} --- you'll see the byte count of
``Hello world!\textbackslash n'' come out as the exit status. This is
the idea I want you to carry forward: \texttt{\%rax} holds the return
value of whatever you just called (syscall or function), and you can
feed it into whatever comes next.

\subsection*{Reading input: \texttt{.bss} and the \texttt{read} syscall}

To read we need somewhere to read \emph{into}, which is what the
\texttt{.bss} section is for --- uninitialized, zero-filled storage
that costs nothing on disk because the loader knows to just hand us
zeroed pages:

\begin{lstlisting}
    .section .bss
buff:   .skip 128
\end{lstlisting}

\texttt{.skip 128} reserves 128 bytes at the label \texttt{buff}.

\texttt{read} is syscall \texttt{0}, and its arguments mirror
\texttt{write}:
\begin{itemize}
  \item \texttt{\%rdi}: fd (0 = stdin)
  \item \texttt{\%rsi}: pointer to the buffer
  \item \texttt{\%rdx}: maximum number of bytes
\end{itemize}
And just like \texttt{write}, \texttt{read} returns the number of
bytes actually read in \texttt{\%rax}.

\subsection*{\texttt{echo.s}: read and then write}

Here's \texttt{echo.s} end-to-end:

\begin{lstlisting}
    .section .bss
buff:   .skip 128

    .section .text
    .global _start

_start:
    mov $0,%rax             # read
    mov $0,%rdi             # stdin
    lea buff(%rip),%rsi
    mov $128,%rdx
    syscall

    mov %rax,%rdx           # bytes-read becomes the count for write
    mov $1,%rax             # write
    lea buff(%rip),%rsi
    mov $1,%rdi             # stdout
    syscall

    xor %rdi,%rdi
    mov $60,%rax
    syscall
\end{lstlisting}

The clever step is the first line after \texttt{syscall}:
\texttt{mov \%rax,\%rdx}. \texttt{read} told us how much it got; we
immediately funnel that into the byte-count argument of \texttt{write},
so we only echo back what was actually typed. If we hardcoded 128
there instead, we'd dump the uninitialized tail of the buffer too ---
not a crash, but garbage.

\subsection*{Jump-based ``subroutines'' --- and why they don't scale}

Before \texttt{call}/\texttt{ret} existed conceptually for us, the
hack was to use labels and unconditional jumps. \texttt{routine.s} is
that hack:

\begin{lstlisting}
rdadder:
    add %rdi,%rdi
    mov %rdi,%rax
    jmp rdret

_start:
    mov $20,%rdi
    jmp rdadder
rdret:
    mov %rax,%rdi
    mov $60,%rax
    syscall
\end{lstlisting}

It works: we jump to \texttt{rdadder}, it doubles \texttt{\%rdi},
then \texttt{jmp rdret} gets us back. But look at what we had to do
--- the ``return address'' is hardcoded into the function body. If we
wanted to call \texttt{rdadder} from two different places, we'd have
to either duplicate it or add some awful switch. It's not reusable,
it's not encapsulated, it's not a function.

\subsection*{Real functions: \texttt{call} and \texttt{ret}}

\texttt{call} and \texttt{ret} are two special instructions that fix
the whole thing. Conceptually:

\begin{itemize}
  \item \texttt{call label} pushes the address of the \emph{next}
    instruction onto the stack and then jumps to \texttt{label}.
  \item \texttt{ret} pops that address off the stack and jumps to it.
\end{itemize}

So \texttt{call}/\texttt{ret} are just ``save your place, jump, come
back'' --- and because the save uses the stack, you can call as many
nested functions as you want.

\texttt{fun1.s} is the first real function we wrote:

\begin{lstlisting}
rdadder:
    add %rdi,%rdi
    mov %rdi,%rax
    ret

_start:
    mov $20,%rdi
    call rdadder         # %rax = 40
    mov %rax,%rdi
    call rdadder         # %rax = 80
    mov %rax,%rdi
    mov $60,%rax
    syscall
\end{lstlisting}

Now \texttt{rdadder} is reusable --- we call it twice in a row and
compound the doubling. Exit status should be 80.

\texttt{fun2.s} is a variant that uses \texttt{\%r9} as a scratch
register instead of clobbering \texttt{\%rdi}:

\begin{lstlisting}
rdadder:
    mov %rdi,%r9
    add %r9,%r9
    mov %r9,%rax
    ret

_start:
    mov $20,%rdi
    call rdadder       # %rdi still 20, %rax = 40
    call rdadder       # %rdi still 20, %rax = 40 again
    mov %rax,%rdi
    mov $60,%rax
    syscall
\end{lstlisting}

Notice what's different from \texttt{fun1.s}: the \texttt{\_start}
here calls \texttt{rdadder} twice back-to-back without reloading
\texttt{\%rdi} in between. Because \texttt{rdadder} no longer
modifies \texttt{\%rdi} (that's the point of using \texttt{\%r9} as a
scratch), both calls see the same input and produce the same output.
Exit status is 40, not 80. If you want to compound the effect, you
have to feed the previous result back in as the next argument ---
exactly what \texttt{fun1.s} does with its \texttt{mov \%rax,\%rdi}
between calls. The lesson: ``the function is pure'' and ``calling it
twice doubles the answer'' are two very different statements.

\subsection*{Local variables mean we need a stack frame}

For anything more interesting than ``double a number,'' functions
need scratch space --- local variables. Every process gets a big
region of memory called the stack, and every function call gets a
contiguous slice of it called a \emph{stack frame}.

Two registers run the show:
\begin{itemize}
  \item \texttt{\%rsp} --- stack pointer, points at the ``top'' of
    the stack.
  \item \texttt{\%rbp} --- base pointer, points at the ``bottom'' of
    the current frame.
\end{itemize}

Important weirdness: the stack \emph{grows downward}. ``Top'' is the
lowest address. Pushing onto the stack \emph{decreases}
\texttt{\%rsp}; popping \emph{increases} it. A function's local
variables live at negative offsets from \texttt{\%rbp}
(e.g.\ \texttt{-8(\%rbp)}, \texttt{-16(\%rbp)}).

\texttt{push} and \texttt{pop} are the two primitives:
\texttt{push \%rbx} subtracts 8 from \texttt{\%rsp} and stores
\texttt{\%rbx} there; \texttt{pop \%rbx} does the inverse.

The canonical function prologue:
\begin{enumerate}
  \item \texttt{push \%rbp} --- save the caller's base pointer.
  \item \texttt{mov \%rsp,\%rbp} --- make our new frame start where
    the stack currently is.
  \item \texttt{sub \$N,\%rsp} --- carve out $N$ bytes of locals.
\end{enumerate}

And the epilogue undoes it:
\begin{enumerate}
  \item \texttt{mov \%rbp,\%rsp} --- drop all the locals in one move.
  \item \texttt{pop \%rbp} --- restore the caller's base pointer.
  \item \texttt{ret}.
\end{enumerate}

\subsection*{\texttt{funStack.s}: a function with locals, by hand}

We're computing $f(x) = 20 + 2x$, but deliberately using local
variables for both \texttt{a = 20} and the intermediate $2x$:

\begin{lstlisting}
fun:
    push %rbp
    mov %rsp,%rbp
    sub $16,%rsp            # room for two quadword locals

    mov %rdi,-8(%rbp)       # local1 = x
    add %rdi,-8(%rbp)       # local1 += x   (so local1 = 2x)
    movq $20,-16(%rbp)      # local2 = 20
    movq -16(%rbp),%rax
    add -8(%rbp),%rax       # %rax = 20 + 2x

    mov %rbp,%rsp
    pop %rbp
    ret

_start:
    mov $10,%rdi
    call fun
    mov %rax,%rdi
    mov $60,%rax
    syscall
\end{lstlisting}

With input 10, this exits with status 40. Walk through it once
on paper: draw a stack growing downward, watch \texttt{push \%rbp},
\texttt{mov \%rsp,\%rbp}, and \texttt{sub \$16,\%rsp} carve out the
frame. The \texttt{-8(\%rbp)} and \texttt{-16(\%rbp)} slots are
\emph{your} two local variables; naming them in a comment alongside
the code is a lifesaver when this gets bigger.

One subtle detail we'll formalize later: after \texttt{push \%rbp}
we've pushed one 8-byte value, which actually re-aligns \texttt{\%rsp}
to a 16-byte boundary (assuming the caller's \texttt{call} left it
8-aligned, which it does). That matters because the ABI requires
\texttt{\%rsp} to be 16-aligned at the point of a \texttt{call},
and some instructions (notably SSE moves) fault if it isn't.

\section{Lecture 6 --- Reusing Functions Across Files, Recursion, and Summing an Array}

With \texttt{call}/\texttt{ret} and stack frames in hand, today was
about doing things that actually resemble programming: linking
pre-written \texttt{readInt}/\texttt{writeInt} helpers in with our
code, writing a recursive \texttt{factorial}, and walking an array to
sum it up.

\subsection*{Logistics}

First, grades for the stack assignments are in --- nice work overall.

Second, the next coding assignment: write an assembly program that
\emph{(i)} declares an array, \emph{(ii)} reads input, \emph{(iii)}
does something interesting combining the two, and \emph{(iv)} exits
cleanly. Two constraints: you must assemble with \texttt{as} and link
with \texttt{ld} (no \texttt{gcc} driver), and you're encouraged to
pull in the \texttt{readInt}/\texttt{writeInt} helpers from the
assembly guide. Due May 10th.

Third, Discussion 3 is on \emph{posits}, an alternative to IEEE-754.
Read \texttt{positTutorial.org} in the repo. The short pitch: posits
spread their precision more evenly across the number line, so they
give up some resolution at very small magnitudes in exchange for
much better resolution at moderate ones, and they ditch the
menagerie of special cases (no $\pm\infty$, no \texttt{NaN}, just a
single \texttt{NaR}). A posit is made up of
\begin{itemize}
  \item a sign bit,
  \item \emph{regime} bits (a kind of super-exponent, variable width),
  \item exponent bits (variable width),
  \item fraction bits (whatever's left).
\end{itemize}
The variable widths are what let posits reallocate precision
dynamically. We'll leave the details for the tutorial.

\subsection*{Linking in functions from other files}

Up to now, everything we've built has been a single \texttt{.s} file
with a \texttt{\_start}. Real programs are many files linked
together. To use a function defined elsewhere you just need to tell
the assembler ``this symbol is defined in another file''; the linker
will wire it up later.

\texttt{readInt.s} defines (and makes global with \texttt{.global})
two functions:

\begin{lstlisting}
    .section .bss
buffer: .zero 128

    .section .text
    .global readInt

parseInt:
    movzbl (%rdi,%rcx,1),%r8d   # read one byte, zero-extended
    sub $'0',%r8
    imul $10,%rax
    add %r8,%rax
    inc %rcx
    cmp %rsi,%rcx
    jb parseInt
    ret

readInt:
    push %rbx
    mov $0,%rax           # sys_read
    mov $0,%rdi           # stdin
    lea buffer(%rip),%rsi
    mov $128,%rdx
    syscall

    mov %rax,%rbx         # bytes read
    mov $0,%rcx
    mov $0,%rax
    lea buffer(%rip),%rdi
    dec %rbx              # trim trailing newline
    mov %rbx,%rsi
    call parseInt
    pop %rbx
    ret
\end{lstlisting}

Two details worth pointing out. First, \texttt{movzbl} (``move byte
to long, zero-extended'') loads a single byte of the input buffer
into the low 8 bits of \texttt{\%r8} and clears the upper bits.
Without the zero extension, garbage from whatever used \texttt{\%r8}
last would corrupt your arithmetic --- classic bug. Second,
\texttt{readInt} does \texttt{push \%rbx} / \texttt{pop \%rbx}
around its body because \texttt{\%rbx} is \emph{callee-saved} in the
Linux ABI. If we clobbered it, the caller would be entitled to
complain.

\texttt{writeInt.s} is the mirror image: it takes an integer in
\texttt{\%rdi}, converts it to decimal digits by repeatedly dividing
by 10, and prints them with \texttt{write}. A few lessons buried in
there:
\begin{itemize}
  \item It allocates a 32-byte buffer \emph{on the stack} and fills
    it \emph{from the end} --- that's how ``print digits in reverse
    order'' naturally handles being emitted left-to-right later.
  \item It handles negatives by remembering a flag and negating
    \texttt{\%rax}, with a special case for \texttt{INT\_MIN}
    ($-2^{63}$), which can't be negated in two's complement. The code
    uses \texttt{movabsq} (not plain \texttt{movq}) because the
    constant $2^{63}$ doesn't fit in a 32-bit sign-extended
    immediate.
  \item It preserves \texttt{\%rbx} (callee-saved again), and uses
    \texttt{leave} as shorthand for ``\texttt{mov \%rbp,\%rsp} then
    \texttt{pop \%rbp}.''
\end{itemize}

To actually call these from another file, we use \texttt{.extern} (a
declaration, not a definition):

\begin{lstlisting}
    .section .text
    .global _start
    .extern readInt
    .extern writeInt

_start:
    call readInt
    mov %rax,%rdi
    call writeInt

    mov $60,%rax
    xor %rdi,%rdi
    syscall
\end{lstlisting}

That's \texttt{readWriteTest.s}. To build it you assemble each file
separately with \texttt{as} and then link them all with \texttt{ld}:

\begin{lstlisting}
as -o readWriteTest.o readWriteTest.s
as -o readInt.o readInt.s
as -o writeInt.o writeInt.s
ld -o readWriteTest readWriteTest.o readInt.o writeInt.o
\end{lstlisting}

Type a number, hit enter, and it prints the same number back. Not
very exciting on its own, but it's the scaffolding for everything
else today.

\subsection*{Recursion: factorial in assembly}

\texttt{fact.s} is the first recursive function we've written. The
idea is familiar --- $n! = n \cdot (n-1)!$, with base case
$0! = 1! = 1$ --- but implementing it in assembly forces you to be
explicit about how the stack carries your intermediate state.

\begin{lstlisting}
fact:
    push %rbp
    mov %rsp,%rbp
    sub $16,%rsp            # one local + padding for alignment

    cmp $1,%rdi
    jle basecase

    mov %rdi,-8(%rbp)       # save n before the recursive call
    dec %rdi                # %rdi = n - 1
    call fact               # recursive call: returns (n-1)! in %rax
    mov -8(%rbp),%rdi       # reload our saved n
    imul %rdi,%rax          # %rax = n * (n-1)!
    jmp leavefact

basecase:
    mov $1,%rax

leavefact:
    leave
    ret
\end{lstlisting}

The essential move is \texttt{mov \%rdi,-8(\%rbp)} before the
recursive call. \texttt{\%rdi} is \emph{caller-saved} --- any
function you call is allowed to clobber it --- and recursion means we
need \texttt{n} again after \texttt{fact(n-1)} returns. So we park it
in our stack frame. Each recursive activation gets its own
\texttt{-8(\%rbp)} slot, because each activation has its own
\texttt{\%rbp}. That's the whole magic trick.

The \texttt{sub \$16,\%rsp} is worth pausing on. We only need 8 bytes
for our one local, but we carve out 16 to keep \texttt{\%rsp} 16-byte
aligned at the next \texttt{call}. If this ever drifts you'll see
mysterious crashes inside library functions that assume SSE
alignment.

The \texttt{\_start} that drives it:

\begin{lstlisting}
_start:
    mov $1,%rax             # print "Enter a number: "
    mov $1,%rdi
    lea message(%rip),%rsi
    mov $16,%rdx
    syscall

    call readInt
    mov %rax,%rdi
    call fact
    mov %rax,%rdi
    call writeInt

    mov $60,%rax
    xor %rdi,%rdi
    syscall
\end{lstlisting}

Try 5 --- you should get 120. Try 20 --- 2432902008176640000. Try 21
--- you'll overflow a 64-bit signed integer and get nonsense. That
silent wraparound is worth remembering.

\subsection*{Summing an array}

\texttt{sumLoopArray.s} pulls together the array addressing from
lecture 4, the counter-loop idiom, and \texttt{writeInt}:

\begin{lstlisting}
    .section .data
arr:    .quad 1,2,3,4

    .section .text
    .global _start
    .extern writeInt

_start:
    mov $0,%rcx             # i = 0
    mov $0,%rax             # accumulator
    lea arr(%rip),%r15      # base pointer

compare:
    cmp $4,%rcx
    jge done
loop:
    add (%r15,%rcx,8),%rax  # %rax += arr[i]
    inc %rcx
    jmp compare
done:
    mov %rax,%rdi
    call writeInt

    mov $60,%rax
    xor %rdi,%rdi
    syscall
\end{lstlisting}

The key line is \texttt{add (\%r15,\%rcx,8),\%rax}: scale-index
addressing again, with scale 8 because we're stepping through
\texttt{.quad} elements. If you change \texttt{arr} to \texttt{.long}
you must change the scale to 4 \emph{and} switch to the 32-bit
forms (\texttt{addl}, \texttt{\%eax}) --- this is the same
size/stride trap I warned about in lecture 4, and it will bite people
in the assignment.

Output for this program is \texttt{10}, because $1+2+3+4=10$. If you
wanted to sum a user-supplied run of numbers instead of a hardcoded
array, you already have all the pieces: \texttt{readInt} in a loop,
accumulate in \texttt{\%rax}, \texttt{writeInt} at the end. That's a
good shape for the upcoming assignment.

\section{Lecture 7 --- Processes, \texttt{fork}, Threads, and Why Sharing State Hurts}

We're pivoting away from assembly for a bit and back into C, but with
a new lens: \emph{concurrency}. Today's arc was processes vs.\ threads,
\texttt{fork} from the bottom up, a quick detour through function
pointers because the threading API needs them, and then a first taste
of why shared state between threads is a minefield.

\subsection*{Logistics}

A few schedule notes worth surfacing: this term, ``finals week'' lands
on what would have been week 10, so I've shifted lectures to fit in
nine weeks. There's no in-class final --- the last week is buffer time
to wrap up the big assignments. Wednesday night of finals week is the
hard cutoff for late homework (concretely: before I sit down to grade
on Thursday morning).

The next major project --- the assembly project from lecture 6 --- is
due in 13 days, on May 10th. Discussions are still due weekly,
quizzes are open and repeatable, and the new course site (still
unpublished) is on GitHub --- check the pinned announcement.

\subsection*{What is a process?}

A \emph{process} is the bundle of context, code, and state that
\emph{is} an executing program: its memory, its open file descriptors,
its registers, its position in the instruction stream, the works.

In the olden days, a process had \emph{a} thread of execution ---
exactly one. Multi-user/multi-tasking systems gave each running
program a single thread and the OS scheduled between them. The OS
\emph{still} schedules: each thread of execution gets a timeslice
before being switched out. What's new (well, new-ish --- this has been
true for decades now) is that one process can host \emph{multiple}
threads of execution, and a real machine has multiple cores that can
literally run threads in parallel.

If you want to peek at how Linux represents all this, the relevant
header is
\href{https://github.com/torvalds/linux/blob/master/include/linux/sched.h}{\texttt{include/linux/sched.h}}.
The interesting fields on a \texttt{task\_struct} for our purposes are
\texttt{parent}, \texttt{thread\_info}, and \texttt{pid}. Linux
internally calls both processes and threads ``tasks'' --- the
distinction at the kernel level is mostly about what they share with
their parent.

\subsection*{\texttt{fork}: the original concurrency primitive}

\texttt{fork} spawns a new process that is a \emph{clone} of the
current one. Same code, same data, same open files, same next
instruction. The implications are stranger than they sound:

\begin{itemize}
  \item The new process picks up at the same line as the parent,
    because it has the same instruction pointer.
  \item There's no shared state between them. Modifying a variable in
    one does not change it in the other. (Modern kernels save space by
    having both processes point at the same physical memory pages
    until somebody writes --- ``copy-on-write'' --- but that's an
    implementation detail.)
  \item It's easy to write \texttt{fork}-style code that doesn't step
    on itself.
  \item It's \emph{harder} to write \texttt{fork}-style code that
    coordinates and shares information, precisely because the state is
    separate.
\end{itemize}

\texttt{fork1.c} is the bare minimum:

\begin{lstlisting}[language=C]
printf("We haven't forked yet\n");
fork();
printf("We've now forked\n");
\end{lstlisting}

Run it and you'll see ``We haven't forked yet'' once and ``We've now
forked'' \emph{twice}. After \texttt{fork()} returns, there are two
processes both sitting at the next line, both about to run that
\texttt{printf}.

\texttt{fork2.c} makes the ``no shared state'' point concrete:

\begin{lstlisting}[language=C]
int num = 5;
fork();
num++;
printf("My number is: %d\n",num);
\end{lstlisting}

Each process has its own copy of \texttt{num}, so each one increments
its own copy and prints \texttt{6}. Not \texttt{7}, not \texttt{6} and
\texttt{7}: each process is bookkeeping its own state, full stop.
(Aside: a polite program would call \texttt{wait()} in the parent so
the child can finish before \texttt{main} returns. We're skipping that
on purpose for now to keep the example short, but it's the kind of
thing that bites you in shell pipelines.)

\subsection*{Process IDs and dispatching}

If \texttt{fork} could only give you two identical clones running in
lockstep, it wouldn't be very useful. The trick that makes it useful
is that \texttt{fork} \emph{returns} different things in the two
processes:

\begin{itemize}
  \item In the parent, \texttt{fork()} returns the child's PID
    (a positive integer).
  \item In the child, \texttt{fork()} returns \texttt{0}.
  \item On failure (e.g.\ out of process slots), it returns
    \texttt{-1} in the original process and there is no child.
\end{itemize}

\texttt{fork3.c} just prints the return value:

\begin{lstlisting}[language=C]
pid_t pid = fork();
printf("The value of pid is: %d\n",pid);
\end{lstlisting}

You'll see two lines: one with \texttt{0} (the child) and one with
some big number (the parent's view of the child's PID).

\texttt{fork4.c} adds \texttt{getpid()}, which always returns ``my own
PID,'' so you can see who's who:

\begin{lstlisting}[language=C]
pid_t pid = fork();
printf("The value of pid is: %d\n",pid);
printf("My name is: %d\n",getpid());
\end{lstlisting}

The child's two lines will say ``pid is 0'' and ``my name is
\texttt{N},'' where \texttt{N} is exactly the number the parent
printed.

\texttt{fork5.c} is the canonical \texttt{fork} idiom: \emph{branch on
the return value} so parent and child do different things:

\begin{lstlisting}[language=C]
pid_t pid = fork();
if(pid != 0){
  printf("I'm the parent!\n");
}
else{
  printf("I'm the child!\n");
}
printf("The value of pid is: %d\n",pid);
printf("My name is: %d\n",getpid());
\end{lstlisting}

This is how every real \texttt{fork}-using program is shaped: the
parent and the child take different branches and do different work.
A shell calls \texttt{fork} to make a child, then has the child
\texttt{exec} the command you typed while the parent waits. A web
server forks one child per incoming connection. The pattern is always
``one call, two return paths.''

\subsection*{For contrast: threads}

Threads are the other half of the picture. A thread of execution
inside an existing process \emph{shares} state with the other threads
in that process: same heap, same globals, same file descriptors. That
makes them lighter to spawn (no copy of the whole address space) and
makes coordination through shared memory trivial --- and it makes
interfering with each other equally trivial. We'll get to that.

The headlines:
\begin{itemize}
  \item Threads share state, processes don't.
  \item Spawning a thread is cheaper than spawning a process.
  \item Threads need explicit control flow: when you make one, you
    have to hand it a function to run, and the thread ends when that
    function returns.
\end{itemize}

That last point is the wedge for the next topic.

\subsection*{A word on function pointers}

To launch a thread you have to give the threading library a
\emph{function} to run. C's way of saying ``here, take this function''
is a function pointer, which a lot of folks haven't seen before.

The point of function pointers is that they let you pass a function as
a first-class argument. If you've used any modern language, you've
probably used ``higher-order functions'' that take other functions ---
\texttt{map}, \texttt{filter}, \texttt{sort}-with-a-comparator,
event handlers. C has had this since forever, just with crunchier
syntax.

\texttt{funpoint1.c} is the warm-up:

\begin{lstlisting}[language=C]
int doubler(int n){
  return 2*n;
}

void maparr(int arr[], int s,
            int (*func)(int)){
  for(int i=0; i<5; i++){
    arr[i] = func(arr[i]);
  }
}

int main(){
  int arr[5] = {0,1,2,3,4};
  maparr(arr,5,doubler);
  for(int i=0; i<5; i++){
    printf("%d\n",arr[i]);
  }
}
\end{lstlisting}

The type \texttt{int (*func)(int)} reads as ``a pointer to a function
that takes an \texttt{int} and returns an \texttt{int}.'' The
parentheses around \texttt{*func} matter --- without them you'd be
declaring a \emph{function returning a pointer to int}, which is a
totally different thing. Inside \texttt{maparr}, \texttt{func(arr[i])}
calls through the pointer; you can also write \texttt{(*func)(arr[i])}
if you want to be explicit, but the syntax sugar lets you drop the
star.

The output is \texttt{0 2 4 6 8}, because \texttt{doubler} got applied
to each element in place.

\subsection*{\texttt{pthread\_create}: actually launching a thread}

Now the threading API. The standard on Linux is POSIX threads
(\texttt{pthreads}). The two calls you'll see today are
\texttt{pthread\_create}, which spawns a thread, and
\texttt{pthread\_join}, which waits for one to finish.

The signature of the thread function is fixed:
\texttt{void* f(void* p)}. It takes a single \texttt{void*} argument
and returns a \texttt{void*}. Why all the \texttt{void*}s? Because
this is C's way of doing polymorphism: any pointer can be cast to
\texttt{void*} and back, so the threading library can hand any
function any kind of data without knowing what either of them is.

\texttt{thread1.c} is the simplest thing that works:

\begin{lstlisting}[language=C]
void* threadFun(void* p){
  printf("Hi from a thread!\n");
  return NULL;
}

int main(){
  pthread_t thread;
  pthread_create(&thread,NULL,threadFun,NULL);
  pthread_join(thread,NULL);
  return 0;
}
\end{lstlisting}

\texttt{pthread\_create} takes:
\begin{enumerate}
  \item A \texttt{pthread\_t*} the library will fill in with the new
    thread's handle.
  \item Optional attributes (\texttt{NULL} = defaults).
  \item The function to run.
  \item The single argument to pass it.
\end{enumerate}
\texttt{pthread\_join} blocks until the thread exits, and optionally
captures its return value (we pass \texttt{NULL} because we don't care).

If you forget the join, \texttt{main} can return before the thread
prints, and depending on timing the program ends without you ever
seeing the message. Always join (or detach) your threads.

You'll also notice the \texttt{return NULL} at the end of
\texttt{threadFun}. Returning a pointer when you have nothing to
return is C's way of doing an ``optional'' value: a non-null pointer
is your result, \texttt{NULL} means ``nothing.'' Modern languages
spell this \texttt{Maybe} or \texttt{Option}; in old-school C it's a
pointer convention.

\subsection*{Passing data into a thread}

\texttt{thread2.c} actually \emph{uses} the argument:

\begin{lstlisting}[language=C]
void* threadFun(void* p){
  int* arg = (int*)p;
  (*arg)++;
  return NULL;
}

int main(){
  int counter = 0;
  pthread_t thread;
  pthread_create(&thread,NULL,threadFun,&counter);
  pthread_join(thread,NULL);
  printf("Counter is: %d\n",counter);
}
\end{lstlisting}

We pass \texttt{\&counter} as the \texttt{void*} argument, the thread
casts it back to \texttt{int*}, increments through it, and the parent
sees \texttt{Counter is: 1} after the join. Crucially, this works
because threads share memory: the thread's \texttt{*arg} and
\texttt{main}'s \texttt{counter} are literally the same bytes. Try
this with \texttt{fork} and you'd see no change in the parent ---
that's the whole difference.

\subsection*{The danger of shared state}

Now the punchline of the lecture. \texttt{thread3.c} launches 100
threads, each of which increments \texttt{counter} 10{,}000 times:

\begin{lstlisting}[language=C]
void* threadFun(void* p){
  int* arg = (int*)p;
  int temp = *arg;
  for(int i=0; i<10000; i++){
    temp++;
  }
  *arg = temp;
  return NULL;
}

int main(){
  int counter = 0;
  pthread_t thread[100];
  for(int i=0; i<100; i++){
    pthread_create(&thread[i],NULL,threadFun,&counter);
  }
  for(int i=0; i<100; i++){
    pthread_join(thread[i],NULL);
  }
  printf("Counter is: %d\n",counter);
}
\end{lstlisting}

You'd hope, naively, for $100 \cdot 10{,}000 = 1{,}000{,}000$. You will
not get $1{,}000{,}000$. Run it a few times and you'll get different
numbers, all far short.

The bug is the read-modify-write pattern in the body. Each thread:
\begin{enumerate}
  \item reads \texttt{*arg} into a local \texttt{temp},
  \item does its 10{,}000 increments on \texttt{temp},
  \item writes \texttt{temp} back to \texttt{*arg}.
\end{enumerate}
While thread A is doing its 10{,}000 increments, threads B, C, D,
\ldots\ are also reading \texttt{*arg}, doing their increments, and
writing back. Whoever writes last wins; everyone else's work is
overwritten. This is a \emph{race condition}, and it is the canonical
failure mode of shared-state concurrency.

\texttt{thread4.c} is the same code with one tiny addition --- a
\texttt{usleep(100)} inside the inner loop, and only 100 iterations
per thread instead of 10{,}000:

\begin{lstlisting}[language=C]
for(int i=0; i<100; i++){
  temp++;
  usleep(100);
}
\end{lstlisting}

The \texttt{usleep} forces the OS to interleave the threads more
visibly --- it's much harder for one thread to ``win'' the race by
finishing its loop in a single timeslice. The result: \texttt{counter}
ends up close to \texttt{100}, not \texttt{10{,}000}. Each of the 100
threads is essentially clobbering everyone else's writes.

The moral isn't ``don't use threads.'' The moral is that \emph{any
time multiple threads touch the same memory, you need some form of
synchronization} --- mutexes, atomics, message passing, something. We
covered the symptom today. The next chunk of the term is going to be
about the cures.
\end{document}
